\documentclass[11pt]{article}

\usepackage{epsfig}
\usepackage{amsfonts}
\usepackage{amssymb}
\usepackage{amstext}
\usepackage{amsmath}
\usepackage{xspace}
\usepackage{theorem}
\usepackage{hyperref}
\usepackage{fullpage}
%\usepackage[]{algorithm2e}
\usepackage{xfrac}
\usepackage{mathtools}
\usepackage{algorithm}
\usepackage{algpseudocode}

\usepackage{listings}

\DeclarePairedDelimiter\ceil{\lceil}{\rceil}
\DeclarePairedDelimiter\floor{\lfloor}{\rfloor}

\usepackage{enumitem}                     

\usepackage{tikz}
\usepackage{tikz-qtree}
\usetikzlibrary{shapes}

\tikzstyle{code} = [black!90, draw=black!30, fill=black!5, very thick,
    rectangle, dashed, inner xsep=10pt, inner ysep=7pt]

\newenvironment{codebox}{
    \hspace{.05\textwidth}
        \begin{tikzpicture}
            \node[code] \bgroup
                \begin{minipage}{.65\textwidth}
                    \begin{alltt}}
                    {\end{alltt}
                \end{minipage}
            \egroup;
        \end{tikzpicture}
}

\newif\ifrubric
\rubrictrue
 
 
 
 \usepackage{titlesec}

\titleformat*{\section}{\bfseries}
\titleformat*{\subsection}{\bfseries}
\titleformat*{\subsubsection}{\bfseries}
\titleformat*{\paragraph}{\bfseries}
\titleformat*{\subparagraph}{\bfseries}


\newenvironment{proof}{{\bf Proof:  }}{\hfill\rule{2mm}{2mm}}
\newenvironment{proofof}[1]{{\bf Proof of #1:  }}{\hfill\rule{2mm}{2mm}}
\newenvironment{proofofnobox}[1]{{\bf#1:  }}{}
\newenvironment{example}{{\bf Example:  }}{\hfill\rule{2mm}{2mm}}

\newtheorem{fact}{Fact}
\newtheorem{lemma}[fact]{Lemma}
\newtheorem{theorem}[fact]{Theorem}
\newtheorem{definition}[fact]{Definition}
\newtheorem{corollary}[fact]{Corollary}
\newtheorem{proposition}[fact]{Proposition}
\newtheorem{claim}[fact]{Claim}
\newtheorem{exercise}[fact]{Exercise}

% math notation
\newcommand{\R}{\ensuremath{\mathbb R}}
\newcommand{\Z}{\ensuremath{\mathbb Z}}
\newcommand{\N}{\ensuremath{\mathbb N}}
\newcommand{\F}{\ensuremath{\mathcal F}}
\newcommand{\SymGrp}{\ensuremath{\mathfrak S}}

\newcommand{\size}[1]{\ensuremath{\left|#1\right|}}
%\newcommand{\ceil}[1]{\ensuremath{\left\lceil#1\right\rceil}}
%\newcommand{\floor}[1]{\ensuremath{\left\lfloor#1\right\rfloor}}
\newcommand{\poly}{\operatorname{poly}}
\newcommand{\polylog}{\operatorname{polylog}}

% anupam's abbreviations
\newcommand{\e}{\epsilon}
\newcommand{\half}{\ensuremath{\frac{1}{2}}}
\newcommand{\junk}[1]{}
\newcommand{\sse}{\subseteq}
\newcommand{\union}{\cup}
\newcommand{\meet}{\wedge}

\newcommand{\prob}[1]{\ensuremath{\text{{\bf Pr}$\left[#1\right]$}}}
\newcommand{\expct}[1]{\ensuremath{\text{{\bf E}$\left[#1\right]$}}}
\newcommand{\Event}{{\mathcal E}}
\newcommand{\E}{\ensuremath{\text{E}}}

\newcommand{\mnote}[1]{\normalmarginpar \marginpar{\tiny #1}}
\newcommand{\opt}[0]{{\text{OPT}}}

\setenumerate[0]{label=(\alph*)}


%%%%%%%%%%%%%%%%%%%%%%%%%%%%%%%%%%%%%%%%%%%%%%%%%%%%%%%%%%%%%%%%%%%%%%%%%%%
% Document begins here %%%%%%%%%%%%%%%%%%%%%%%%%%%%%%%%%%%%%%%%%%%%%%%%%%%%
%%%%%%%%%%%%%%%%%%%%%%%%%%%%%%%%%%%%%%%%%%%%%%%%%%%%%%%%%%%%%%%%%%%%%%%%%%%



\begin{document}

\noindent {\large {\bf 601.433/633 Introduction to Algorithms} \hfill {{\bf Fall 2026}}}\\ 
{{\bf Homework \#7}} \hfill {{\bf Due:} November 5, 2026, 11:59pm} \\
\rule[0.1in]{\textwidth}{0.4pt}

Remember: this is an oral presentation homework!  You are required to work in groups of three, and may collaborate only with other students in your group.  We encourage you to write up solutions, but your grade will be based entirely on your presentation.  At your presentation, the grader will randomly assign each student to present a problem.  You should be able to create a group and sign up for a slot using Canvas.  The deadline for signing up for a presentation slot is 5pm Wed Nov 4.

You may use the internet or LLMs to look up formulas, definitions, concepts, etc., but may not simply look up the answers online or ask an LLM.  

\noindent \rule[0.1in]{\textwidth}{0.4pt}

\section{Faster Shortest Paths}

Let $G = (V, E)$ be a directed graph with lengths $\ell : E \rightarrow \R$ with no negative-length cycles.  Let $v \in V$, and let $\alpha$ denote the maximum, over all $u \in V$, of the number of edges on a shortest path from $v$ to $u$ (where the shortest-path is defined with respect to the weights).  Given $G$, $w$, and $v$ (but not $\alpha$), give an algorithm that computes shortest path distances from $v$ to all other nodes in $O(m \alpha)$ time.  You may assume that $m = \Omega(n)$.  Prove correctness and running time.



\section{Multi-Set Shortest Paths}
Let $G = (V, E)$ be a directed graph with weighted edges; edge weights can be positive, negative, or zero.  There are no negative-weight cycles.  Suppose the vertices of $G$ are partitioned into $k$ disjoint subsets $V_1, V_2, \dots, V_k$; that is, every vertex of $G$ belongs to exactly one subset $V_i$. For each $i, j \in [k]$, let $\delta(i, j)$ denote the minimum shortest-path distance between vertices in $V_i$ and vertices in $V_j$:
\[
\delta(i,j) = \min \left\{d(v_i, v_j) \mid v_i \in V_i\ \text{and}\ v_j \in V_j\right\}
\]

Design an algorithm to compute $\delta(i,j)$ for all $i,j \in [k]$ that runs in time $O(mn + kn\log n)$.  You may assume that $m = \Omega(n)$.  Prove correctness and running time.


\section{Single-Source Shortest Paths}
Let $G = (V, E)$ be a connected, undirected graph with positive edge lengths $\ell : E \rightarrow \mathbb{R}^+$.  This graph represents a road network.  You are trying to get from $s \in V$ to $t \in V$.  Unfortunately, you are a coffee addict: getting caffeine often enough is the most important thing to you.  In fact, there is some number $L \in \mathbb{R}^+$ so that if you go more that $L$ distance without getting coffee, you will get a caffeine migraine and crash your car.  So you simply \emph{cannot} go more than $L$ distance without stopping for coffee.  Fortunately, there is a coffee shop at every vertex $v \in V$ (including $s$ and $t$) where you can get coffee.  However, you are also a coffee snob, and you \emph{hate} Starbucks.  And there is a special subset of nodes $S \subseteq V$ where the only coffee shop on that node is a Starbucks.

So you have the following problem: given a subset $S \subseteq V$ of vertices where the only coffee shop is a Starbucks, find the path from $s$ to $t$ which uses only edges of length at most $L$, and has the \emph{fewest} vertices in $S$.  You may assume that $s,t\not\in S$.  Design an $O(m + n\log n)$-time algorithm for this problem, and prove running time and correctness.




\end{document}

